\worksheettitle
  {Составление чисел из цифр}
  {\modulemark{M05} Учебный модуль}

\studentnameblock

\begin{notebox}[Чему вы научитесь]
  Точно читать условие задачи, строить полный упорядоченный список чисел,
  учитывать повторения и обязательные цифры, а также находить пропуски и
  повторы в готовом списке.
\end{notebox}

\begin{checkpointbox}[Вход]
  В числе \(525\) подчеркните цифру десятков. Её значение:
  \answerblank[20mm].

  Почему числа \(25\) и \(52\) различны?
  \answerblank[82mm].
\end{checkpointbox}

\begin{notebox}[Как устроен маршрут]
  \textbf{Обязательное ядро:} вход, разобранные примеры, алгоритм, практика с
  опорой, самостоятельные задания, разбор ошибок и контрольная точка.

  \textbf{Дополнительная практика:} задания, отмеченные этим названием. Их
  выполняйте по указанию учителя или если нужно закрепить способ.
\end{notebox}

\section{Сначала уточняем условие}

\begin{problembox}[Разобранный пример: один набор цифр, разные задачи]
  Используем только цифры \(2\) и \(5\).
  \begin{center}
  \renewcommand{\arraystretch}{1.35}
  \begin{tabular}{>{\raggedright\arraybackslash}p{54mm}|p{83mm}}
    \toprule
    условие & результат \\
    \midrule
    Двузначные числа; повторения разрешены; использовать обе цифры
    необязательно. & \(22,\ 25,\ 52,\ 55\) \\
    Двузначные числа; повторения запрещены. & \(25,\ 52\) \\
    Трёхзначные числа; повторения разрешены; обе цифры обязательны, то есть
    каждая должна встретиться хотя бы один раз. &
    \(225,\ 252,\ 255,\ 522,\ 525,\ 552\) \\
    \bottomrule
  \end{tabular}
  \end{center}
\end{problembox}

\begin{notebox}[Что сделали и что нужно проверить]
  \begin{enumerate}[label=\textbf{\arabic*.}]
    \item Какие цифры разрешены?
    \item Сколько цифр должно быть в числе?
    \item Разрешены ли повторения?
    \item Обязательно ли использовать каждую заданную цифру?
  \end{enumerate}
\end{notebox}

\newpage
\section{Строим полный упорядоченный список}

\begin{reflectionbox}[Правило: алгоритм полного списка]
  \begin{enumerate}[label=\textbf{\arabic*.}]
    \item Выпишите разрешённые первые цифры в порядке возрастания.
    \item Зафиксируйте первую цифру.
    \item Переберите оставшиеся позиции по порядку. При запрете повторений
      исключайте уже использованную цифру.
    \item Перейдите к следующей допустимой первой цифре.
    \item Проверьте количество вариантов, отсутствие пропусков и повторов.
  \end{enumerate}
\end{reflectionbox}

\begin{problembox}[Пример без повторений]
  Из цифр \(2,\ 5,\ 7\) составим все трёхзначные числа, используя каждую цифру
  ровно один раз.
  \begin{center}
  \renewcommand{\arraystretch}{1.5}
  \begin{tabular}{c|c|c}
    \toprule
    первая цифра 2 & первая цифра 5 & первая цифра 7 \\
    \midrule
    \(257\) & \(527\) & \(725\) \\
    \(275\) & \(572\) & \(752\) \\
    \bottomrule
  \end{tabular}
  \end{center}
  Для первой позиции есть три выбора. После каждого выбора оставшиеся цифры
  можно расположить двумя способами. Получено \(3\cdot2=6\) чисел.
\end{problembox}

\begin{practicebox}[1. Практика с опорой: выполните по алгоритму]
  Из цифр \(1,\ 3,\ 6\) составьте все двузначные числа без повторений.

  Начинаются с 1: \answerblank[38mm].

  Начинаются с 3: \answerblank[38mm].

  Начинаются с 6: \answerblank[38mm].

  Полный список: \answerblank[86mm]. Всего чисел: \answerblank[14mm].
\end{practicebox}

\begin{practicebox}[2. Объясните полноту]
  Почему в предыдущем списке нет пропусков и повторов?
  \writinglines{2}
\end{practicebox}

\newpage
\section{Самостоятельная практика: учитываем ноль}

\begin{warningbox}[Почему ветвь с нулём исключаем]
  Записи \(038\) и \(083\) обозначают числа \(38\) и \(83\). Поэтому они не
  являются записями трёхзначных чисел. Запись многозначного натурального числа
  не начинают с нуля.
\end{warningbox}

\begin{problembox}[Пример]
  Из цифр \(0,\ 3,\ 8\) составим все трёхзначные числа без повторений,
  используя каждую цифру ровно один раз.
  \begin{center}
  \renewcommand{\arraystretch}{1.45}
  \begin{tabular}{c|c|c}
    \toprule
    первая цифра 0 & первая цифра 3 & первая цифра 8 \\
    \midrule
    недопустимо & \(308\) & \(803\) \\
    недопустимо & \(380\) & \(830\) \\
    \bottomrule
  \end{tabular}
  \end{center}
  Полный список: \(308,\ 380,\ 803,\ 830\).
\end{problembox}

\begin{practicebox}[1. Составьте списки]
  \begin{enumerate}[label=\textbf{\arabic*.}]
    \item Все двузначные числа из цифр \(1\) и \(6\); повторения разрешены,
      использовать обе цифры необязательно: \answerblank[72mm].
    \item Все двузначные числа из цифр \(1,\ 6,\ 9\) без повторений:
      \answerblank[86mm].
    \item Все трёхзначные числа из цифр \(0,\ 4,\ 9\), используя каждую ровно
      один раз: \answerblank[86mm].
  \end{enumerate}
\end{practicebox}

\begin{warningbox}[2. Найдите причину ошибки]
  Для цифр \(3\) и \(7\), когда повторения разрешены и использовать обе цифры
  необязательно, записали только \(37\) и \(73\).

  Потерянные числа: \answerblank[42mm].
  Причина: \answerblank[74mm].
\end{warningbox}

\begin{practicebox}[3. Обе цифры обязательны]
  Из цифр \(2\) и \(9\) составьте все трёхзначные числа. Повторения разрешены,
  обе цифры обязательны:
  \answerblank[100mm].
\end{practicebox}

\newpage
\section{Разбираемся в ошибках и сравниваем списки}

\begin{practicebox}[Дополнительная практика: один набор, два условия]
  Используйте цифры \(1,\ 4,\ 7\).
  \begin{enumerate}[label=\textbf{\arabic*.}]
    \item Двузначные числа без повторений. Ожидаемое количество:
      \answerblank[13mm]. Список: \answerblank[82mm].
    \item Двузначные числа с повторениями; использовать все цифры необязательно.
      Ожидаемое количество: \answerblank[13mm]. Список:
      \answerblank[82mm].
    \item Какие числа добавились и почему?
      \writinglines{1}
  \end{enumerate}
\end{practicebox}

\begin{warningbox}[2. Найдите причину и исправьте список]
  Из цифр \(1,\ 3,\ 6\) составляли двузначные числа без повторений:
  \[
    13,\ 16,\ 31,\ 36,\ 61,\ 61.
  \]
  Повтор: \answerblank[15mm]. Пропущено: \answerblank[15mm].
  Как структура списка помогает это увидеть?
  \writinglines{1}
\end{warningbox}

\begin{warningbox}[3. Назовите причину и исключите лишние записи]
  Для цифр \(0,\ 2,\ 6\) без повторений записали:
  \(02,\ 06,\ 20,\ 26,\ 60,\ 62\).

  Исключить: \answerblank[30mm]. Останется чисел: \answerblank[13mm].
  Почему исправленный список полный?
  \writinglines{1}
\end{warningbox}

\begin{practicebox}[Дополнительная практика: составьте условие по ответу]
  Для списка \(33,\ 35,\ 53,\ 55\) придумайте точное условие:
  \writinglines{1}
  Запишите новый список, если повторения запретить: \answerblank[48mm].
\end{practicebox}

\newpage
\section{Контрольная точка}

\begin{checkpointbox}[Выполните без подсказки]
  \begin{enumerate}[label=\textbf{\arabic*.}]
    \item Из цифр \(1\) и \(4\) составьте все двузначные числа. Повторения
      разрешены, использовать обе цифры необязательно:
      \answerblank[70mm].
    \item Из цифр \(0,\ 4,\ 7\) составьте все трёхзначные числа без повторений,
      используя каждую цифру ровно один раз:
      \answerblank[82mm].
    \item Из цифр \(2\) и \(9\) составьте все трёхзначные числа с повторениями,
      если обе цифры обязательны:
      \answerblank[92mm].
    \item В списке \(25,\ 27,\ 52,\ 57,\ 72,\ 72\) найдите повтор и пропуск,
      если использованы цифры \(2,\ 5,\ 7\) без повторений.
      \writinglines{1}
    \item Объясните, почему списки в пунктах с 1-го по 3-й полны, почему в пункте 2 нет
      числа, начинающегося с нуля, и как списки упорядочены.
      \writinglines{3}
  \end{enumerate}
\end{checkpointbox}

\begin{reflectionbox}[Проверяю себя]
  \choicebox\ Уточняю условие о повторениях и обязательных цифрах.

  \choicebox\ Группирую числа по первой цифре и соблюдаю порядок.

  \choicebox\ Не начинаю запись числа с нуля.

  \choicebox\ Нахожу пропуски и повторы по структуре списка.
\end{reflectionbox}

\begin{notebox}[Дальнейший маршрут]
  Если условие изменялось, список строился случайно, появился ведущий ноль
  или остались пропуски и повторы, выполните M05-R. Если способ понятен, но
  нужна тренировка, выполните дополнительную практику или M05-H. Если
  контроль выполнен устойчиво и полнота объяснена, переходите к M06;
  M05\(\star\) можно выполнить дополнительно.
\end{notebox}
